% !TEX program = xelatex
\documentclass[a4paper,10pt]{article}
\usepackage[UTF8,fontset=none]{ctex}
\usepackage{xeCJK}
\setCJKmainfont{Songti SC}
\setCJKsansfont{STHeiti}
\setCJKmonofont{Songti SC}
\usepackage[margin=2.2cm]{geometry}
\usepackage{tikz}
\usetikzlibrary{shapes.geometric, arrows.meta, positioning, calc, fit}
\usepackage{amsmath, amssymb}
\usepackage{longtable}
\usepackage{array}
\usepackage{booktabs}
\usepackage{multirow}
\usepackage{enumitem}
\usepackage{fancyhdr}
\usepackage{titlesec}
\usepackage{hyperref}
\usepackage{xcolor}
\usepackage{tabularx}

\hypersetup{
  colorlinks=true,
  linkcolor=blue!60!black,
  urlcolor=blue!60!black,
  bookmarksnumbered=true
}

% 页眉页脚
\pagestyle{fancy}
\fancyhf{}
\fancyhead[L]{\small normalizeAngle~---~函数设计文档}
\fancyhead[R]{\small 内部公开}
\fancyfoot[C]{\thepage}
\renewcommand{\headrulewidth}{0.4pt}

% 层级编号
\setcounter{secnumdepth}{3}
\setcounter{tocdepth}{3}

% 表格列宽辅助
\newcolumntype{L}[1]{>{\raggedright\arraybackslash}p{#1}}
\newcolumntype{C}[1]{>{\centering\arraybackslash}p{#1}}

% ============================================================
\begin{document}

% ==================== 封面表格 ====================
\thispagestyle{empty}

\begin{center}
{\CJKfamily{zhsong}\bfseries\Large 函数设计文档}
\end{center}

\vspace{8pt}

\begin{center}
\renewcommand{\arraystretch}{1.5}
\begin{tabular}{|L{3.2cm}|L{4.5cm}|L{3.2cm}|L{4.5cm}|}
\hline
\textbf{函数英文名} & normalizeAngle &
\textbf{函数中文名} & 角度归一化 \\
\hline
\textbf{函数类型} & 元件 (Element) &
\textbf{版本} & 2.0 \\
\hline
\textbf{作者} & 许庆 &
\textbf{日期} & 2026-06-26 \\
\hline
\textbf{联系方式} & \multicolumn{3}{l|}{18612426814~/~qingxu@tsinghua.edu.cn} \\
\hline
\textbf{密级} & 内部公开 &
\textbf{AI辅助} & Cursor Agent (Claude Opus 4.6) \\
\hline
\end{tabular}
\end{center}

\vspace{12pt}

\begin{center}
\renewcommand{\arraystretch}{1.4}
\begin{tabular}{|C{1.2cm}|C{2cm}|L{5cm}|C{2cm}|C{2cm}|}
\hline
\multicolumn{5}{|c|}{\textbf{更改记录}} \\
\hline
\textbf{版本} & \textbf{日期} & \textbf{更改说明} & \textbf{作者} & \textbf{辅助AI} \\
\hline
2.0 & 2026-06-26 & 初始版本 & 许庆 & Claude Opus 4.6 \\
\hline
\end{tabular}
\end{center}

\newpage

% ==================== 目录 ====================
\tableofcontents
\newpage

% ============================================================
%  1. 函数需求
% ============================================================
\section{函数需求}

\subsection{功能概述}

\texttt{normalizeAngle} 将任意角度值归一化到 $[-\pi, \pi]$ 区间。
利用 $\text{atan2}(\sin\theta, \cos\theta)$ 的数学性质实现角度的周期折叠。

\subsection{输入/输出/返回值}

\renewcommand{\arraystretch}{1.3}
\begin{longtable}{|C{1.5cm}|L{3.5cm}|C{2cm}|L{8cm}|}
\hline
\textbf{类别} & \textbf{名称} & \textbf{类型} & \textbf{说明} \\
\hline
\endfirsthead
\hline
\textbf{类别} & \textbf{名称} & \textbf{类型} & \textbf{说明} \\
\hline
\endhead
输入 & \texttt{angle} & double & 待归一化角度 (rad)，取值范围任意 \\
\hline
返回值 & --- & double & 归一化后的角度 (rad)，$\in [-\pi, \pi]$ \\
\hline
\end{longtable}

% ============================================================
%  1.2 算法设计
% ============================================================
\subsection{算法设计}

角度归一化的核心公式为：
\begin{equation}
  \text{result} = \text{atan2}\!\big(\sin(\theta),\; \cos(\theta)\big)
\end{equation}

\textbf{数学原理}：对于任意角度 $\theta$，$(\cos\theta, \sin\theta)$ 是单位圆上的一点，
$\text{atan2}$ 函数返回该点对应的角度，值域恒为 $(-\pi, \pi]$。
这等价于将角度模 $2\pi$ 后映射到 $[-\pi, \pi]$：
\begin{equation}
  \text{result} \equiv \theta \pmod{2\pi}, \quad \text{result} \in [-\pi, \pi]
\end{equation}

该方法的优点：
\begin{itemize}
  \item 无需循环或条件分支，单次三角函数调用完成
  \item 对大角度值（多圈）同样有效
  \item 自动处理符号和象限
\end{itemize}

% ============================================================
%  1.3 程序流程图
% ============================================================
\subsection{程序流程图}

\begin{center}
\begin{tikzpicture}[
  node distance=1.0cm,
  startstop/.style={rectangle, rounded corners, minimum width=3.5cm, minimum height=0.7cm,
    text centered, draw=black, fill=blue!8, font=\small},
  process/.style={rectangle, minimum width=5cm, minimum height=0.7cm,
    text centered, draw=black, fill=orange!8, font=\small},
  arrow/.style={-{Stealth[length=2.5mm]}, thick}
]

\node (start) [startstop] {开始};
\node (input) [process, below=of start] {读入 angle};
\node (calc) [process, below=of input] {result = atan2(sin(angle), cos(angle))};
\node (output) [process, below=of calc] {返回 result};
\node (stop) [startstop, below=of output] {结束};

\draw [arrow] (start) -- (input);
\draw [arrow] (input) -- (calc);
\draw [arrow] (calc) -- (output);
\draw [arrow] (output) -- (stop);

\end{tikzpicture}
\end{center}

% ============================================================
%  1.4 函数接口
% ============================================================
\subsection{函数接口}

\begin{verbatim}
double normalizeAngle(const double angle);
\end{verbatim}

% ============================================================
%  1.5 输入输出模块图
% ============================================================
\subsection{输入输出模块图}

\begin{center}
\begin{tikzpicture}[
  box/.style={rectangle, draw=black, thick, minimum width=3.0cm, minimum height=0.6cm,
    text centered, fill=yellow!10, font=\small},
  core/.style={rectangle, draw=black, very thick, minimum width=4.5cm, minimum height=1.5cm,
    text centered, fill=orange!10, font=\normalsize\bfseries, rounded corners=3pt},
  arrow/.style={-{Stealth[length=2.5mm]}, thick}
]

\node (core) [core] {normalizeAngle};

\node (in1) [box, left=3cm of core] {angle (rad)};

\node (out1) [box, right=3cm of core] {result (rad)};

\draw [arrow] (in1) -- (core);
\draw [arrow] (core) -- (out1);

\node [above=0.1cm of in1, font=\small\bfseries] {输入 (Input)};
\node [above=0.1cm of out1, font=\small\bfseries] {输出 (Output)};

\end{tikzpicture}
\end{center}

% ============================================================
%  1.6 函数诸元表
% ============================================================
\subsection{函数诸元表}

\renewcommand{\arraystretch}{1.3}
\begin{longtable}{|L{3.8cm}|L{11.5cm}|}
\hline
\textbf{属性} & \textbf{说明} \\
\hline
\endfirsthead
\hline
\textbf{属性} & \textbf{说明} \\
\hline
\endhead
函数英文名 & \texttt{normalizeAngle} \\
\hline
函数中文名 & 角度归一化 \\
\hline
函数类型 & 元件 (Element) \\
\hline
版本 & 2.0 \\
\hline
源文件 & \texttt{src/cpp/normalizeAngle/normalizeAngle.cpp} \\
\hline
头文件 & \texttt{include/cpp/normalizeAngle/normalizeAngle.hpp} \\
\hline
命名空间 & \texttt{waypoint\_follow} \\
\hline
输入数量 & 1（angle） \\
\hline
输出数量 & 1（返回值 double） \\
\hline
参数数量 & 0 \\
\hline
状态数量 & 0（无状态，纯函数） \\
\hline
下级函数数量 & 0 \\
\hline
动态内存分配 & 无 \\
\hline
外部依赖 & \texttt{<cmath>}（std::sin, std::cos, std::atan2） \\
\hline
算法类别 & 角度处理 / 三角函数映射 \\
\hline
\end{longtable}

% ============================================================
%  1.7 函数架构
% ============================================================
\subsection{函数架构}

无下级函数调用。\texttt{normalizeAngle} 为叶节点元件，仅调用标准库三角函数。

% ============================================================
%  1.8 数据结构定义
% ============================================================
\subsection{数据结构定义}

无自定义数据结构。输入输出均为 \texttt{double} 标量。

% ============================================================
%  1.9 测试用例
% ============================================================
\subsection{测试用例}

\subsubsection{正常测试用例}

{\small
\begin{longtable}{|C{0.6cm}|L{2.2cm}|L{4.5cm}|L{3.5cm}|L{3.8cm}|}
\hline
\textbf{编号} & \textbf{场景} & \textbf{输入} & \textbf{预期输出} & \textbf{验证准则} \\
\hline
\endfirsthead
\hline
\textbf{编号} & \textbf{场景} & \textbf{输入} & \textbf{预期输出} & \textbf{验证准则} \\
\hline
\endhead
N-1 &
零角度 &
$\theta{=}0$ &
$r = 0$ &
$|r| < \epsilon$ \\
\hline
N-2 &
正角度（区间内） &
$\theta{=}\pi/4$ &
$r = \pi/4$ &
$|r - \pi/4| < \epsilon$ \\
\hline
N-3 &
负角度（区间内） &
$\theta{=}{-}\pi/2$ &
$r = -\pi/2$ &
$|r + \pi/2| < \epsilon$ \\
\hline
N-4 &
超出上限（$3\pi/2$） &
$\theta{=}3\pi/2$ &
$r = -\pi/2$ &
$|r + \pi/2| < \epsilon$ \\
\hline
N-5 &
多圈正角度 &
$\theta{=}5\pi$ &
$r = \pi$（或 $-\pi$） &
$|r| - \pi < \epsilon$ \\
\hline
\end{longtable}
}

\vspace{6pt}
{\small
\begin{longtable}{|L{2.5cm}|L{13cm}|}
\hline
\multicolumn{2}{|c|}{\textbf{正常测试用例符号说明}} \\
\hline
\textbf{符号} & \textbf{含义} \\
\hline
$\theta$ & 输入角度 angle (rad) \\
\hline
$r$ & 返回值 result (rad) \\
\hline
$\epsilon$ & 浮点误差容限，取 $10^{-12}$ \\
\hline
$\pi$ & 圆周率 $\approx 3.14159265358979$ \\
\hline
\end{longtable}
}

\subsubsection{异常测试用例}

{\small
\begin{longtable}{|C{0.6cm}|L{2.0cm}|L{4.2cm}|L{4.0cm}|L{3.8cm}|}
\hline
\textbf{编号} & \textbf{场景} & \textbf{输入} & \textbf{预期输出/行为} & \textbf{验证准则} \\
\hline
\endfirsthead
\hline
\textbf{编号} & \textbf{场景} & \textbf{输入} & \textbf{预期输出/行为} & \textbf{验证准则} \\
\hline
\endhead
A-1 &
NaN 输入 &
$\theta{=}\text{NaN}$ &
$r = \text{NaN}$ &
sin/cos(NaN) = NaN, atan2(NaN,NaN) = NaN \\
\hline
A-2 &
$+\infty$ 输入 &
$\theta{=}+\infty$ &
$r = \text{NaN}$ &
sin($\infty$) 为 NaN（C标准） \\
\hline
A-3 &
$-\infty$ 输入 &
$\theta{=}-\infty$ &
$r = \text{NaN}$ &
sin($-\infty$) 为 NaN（C标准） \\
\hline
A-4 &
极大正值 &
$\theta{=}1.0\!\times\!10^{15}$ &
$r \in [-\pi, \pi]$ &
结果在有效范围内 \\
\hline
A-5 &
极大负值 &
$\theta{=}{-}1.0\!\times\!10^{15}$ &
$r \in [-\pi, \pi]$ &
结果在有效范围内 \\
\hline
A-6 &
精确 $\pi$ 边界 &
$\theta{=}\pi$ &
$r = \pi$ &
$|r - \pi| < \epsilon$ \\
\hline
A-7 &
精确 $-\pi$ 边界 &
$\theta{=}{-}\pi$ &
$r = -\pi$（或 $\pi$） &
$|r| - \pi < \epsilon$ \\
\hline
A-8 &
略超 $\pi$ &
$\theta{=}\pi + 10^{-14}$ &
$r \approx -\pi$（翻转） &
$|r + \pi| < 10^{-10}$ \\
\hline
A-9 &
subnormal 值 &
$\theta{=}5\!\times\!10^{-324}$ &
$r \approx 0$ &
$|r| < \epsilon$ \\
\hline
A-10 &
接近 DBL\_MAX &
$\theta{=}1.7\!\times\!10^{308}$ &
$r \in [-\pi, \pi]$ 或 NaN &
函数不崩溃 \\
\hline
\end{longtable}
}

\vspace{6pt}
{\small
\begin{longtable}{|L{2.5cm}|L{13cm}|}
\hline
\multicolumn{2}{|c|}{\textbf{异常测试用例符号说明}} \\
\hline
\textbf{符号} & \textbf{含义} \\
\hline
$\theta$ & 输入角度 angle (rad) \\
\hline
$r$ & 返回值 result (rad) \\
\hline
NaN & IEEE 754 非数值 (Not a Number) \\
\hline
$+\infty$/$-\infty$ & IEEE 754 正/负无穷大 \\
\hline
$\epsilon$ & 浮点误差容限，取 $10^{-12}$ \\
\hline
subnormal & 非规格化浮点数（极接近零） \\
\hline
DBL\_MAX & double 最大有限值 $\approx 1.7\!\times\!10^{308}$ \\
\hline
\end{longtable}
}

\end{document}
